\section{Modelo matemático}

\noindent En esta sección se presenta un modelo matemático para el Multi-Source Capacitated Facility Location Problem with Customer Incompatibilities (MS-CFLP-CI). Este modelo representa las decisiones de apertura de bodegas, asignación de demanda de clientes y control de incompatibilidades. La formulación corresponde a un modelo de programación lineal entera mixta, ya que combina variables binarias con variables de cantidad o flujo \cite{maia2023metaheuristic,mansini2022kernel,ceschia2024multineighborhood}.

\subsection{Conjuntos}

\noindent Sea $I$ el conjunto de bodegas candidatas y sea $J$ el conjunto de clientes. Además, se define $P$ como el conjunto de pares de clientes incompatibles. Si $(a,b) \in P$, entonces los clientes $a$ y $b$ no pueden ser atendidos por una misma bodega.

\subsection{Parámetros}

\noindent Los parámetros considerados son los siguientes:

\begin{itemize}
    \item $C_i$: capacidad máxima de la bodega $i$.
    \item $FC_i$: costo fijo de apertura de la bodega $i$.
    \item $G_j$: demanda total del cliente $j$.
    \item $SC_{ji}$: costo unitario de operación por satisfacer demanda del cliente $j$ desde la bodega $i$.
\end{itemize}

\subsection{Variables de decisión}

\noindent El modelo utiliza las siguientes variables:

\begin{itemize}
    \item $y_i \in \{0,1\}$: toma valor $1$ si la bodega $i$ es abierta, y $0$ en caso contrario.
    \item $x_{ji} \in \mathbb{Z}_{\geq 0}$: cantidad de demanda del cliente $j$ satisfecha desde la bodega $i$.
    \item $z_{ji} \in \{0,1\}$: toma valor $1$ si el cliente $j$ recibe una cantidad positiva de productos desde la bodega $i$, y $0$ en caso contrario.
\end{itemize}

\noindent La variable $x_{ji}$ permite representar la característica principal del MS-CFLP-CI, ya que un cliente puede recibir demanda desde más de una bodega. Por otra parte, la variable $z_{ji}$ permite identificar si existe una asignación efectiva entre un cliente y una bodega, lo que resulta necesario para controlar las incompatibilidades entre clientes.

\subsection{Función objetivo}

\noindent El objetivo es minimizar el costo total de la solución, compuesto por los costos fijos de apertura de bodegas y los costos operacionales asociados al envío de demanda.

\begin{equation}
\min
\left\{
\sum_{i \in I} FC_i y_i
+
\sum_{j \in J} \sum_{i \in I} SC_{ji} x_{ji}
\right\}
\end{equation}

\noindent El primer término representa el costo de abrir las bodegas utilizadas, mientras que el segundo término representa el costo de satisfacer la demanda de los clientes desde dichas bodegas.

\subsection{Restricciones}

\noindent Cada cliente debe tener su demanda completamente satisfecha. La demanda puede ser cubierta por una o varias bodegas.

\begin{equation}
\sum_{i \in I} x_{ji} = G_j \quad \forall j \in J
\end{equation}

\noindent Cada bodega solo puede operar hasta su capacidad máxima. Además, si una bodega no está abierta, no puede enviar productos a ningún cliente.

\begin{equation}
\sum_{j \in J} x_{ji} \leq C_i y_i \quad \forall i \in I
\end{equation}

\newpage

\noindent Las siguientes restricciones relacionan la cantidad enviada con la variable de asignación. Si el cliente $j$ no está asignado a la bodega $i$, entonces no puede recibir demanda desde ella. A su vez, si $z_{ji}$ toma valor $1$, se exige que exista una cantidad positiva enviada desde la bodega $i$ hacia el cliente $j$.

\begin{equation}
x_{ji} \leq G_j z_{ji} \quad \forall j \in J, \forall i \in I
\end{equation}

\begin{equation}
x_{ji} \geq z_{ji} \quad \forall j \in J, \forall i \in I
\end{equation}

\noindent También se exige que un cliente solo pueda ser atendido por una bodega que esté abierta.

\begin{equation}
z_{ji} \leq y_i \quad \forall j \in J, \forall i \in I
\end{equation}

\noindent Finalmente, se incorporan las incompatibilidades entre clientes. Si dos clientes son incompatibles, no pueden ser atendidos por una misma bodega.

\begin{equation}
z_{ai} + z_{bi} \leq 1 \quad \forall (a,b) \in P, \forall i \in I
\end{equation}

\noindent El dominio de las variables queda definido como:

\begin{equation}
y_i \in \{0,1\} \quad \forall i \in I
\end{equation}

\begin{equation}
z_{ji} \in \{0,1\} \quad \forall j \in J, \forall i \in I
\end{equation}

\begin{equation}
x_{ji} \in \mathbb{Z}_{\geq 0} \quad \forall j \in J, \forall i \in I
\end{equation}
